%% kapa: exam=1 round=19 year=2008 period=1 form=A subject=appraisal_law subject_ko=부동산관계법규 seq=26 number=66
%% source: Q-Net 감정평가사 기출문제 (cst00309) artlSeq=1001962 file=extracted/1008/1001962/2008년 제19회 감정평가사 1차시험 기출문제/1교시.hwp
%% conversion: kordoc 4.24.51 -> tools/gen1.py (auto); review=proofread-sonnet; flags=none
\begin{question}{66}
국유재산법령상 옳지 않은 것은? (다툼이 있는 경우에는 판례에 의함)
\begin{choices}
  \item 국유재산 무단점유자에 대한 변상금부과처분에 있어서 그 납부고지서 또는 사전통지서에 변상금 산출근거를 명시하지 않았더라도 적법한 처분이 된다.
  \item 관리청은 행정재산 또는 보존재산의 관리를 위하여 필요한 경우에는 국가기관외의 자에게 당해 재산의 관리를 위탁할 수 있다.
  \item 토지의 지목이 도로이고 국유재산대장에 등재되어 있다는 사정만으로 바로 그 토지가 도로로서 행정재산에 해당한다고 할 수는 없다.
  \item 은닉국유재산으로 의심되는 토지의 목록을 제출하면서 구두로 그에 대한 전반적인 조사를 요청한 것이 단서가 되어 해당 토지에 대한 국가 환수가 이루어졌다면 위 구두에 의한 신고행위는 국유재산법 제53조(은닉재산등의 신고)에서 규정하는 적법한 신고로 볼 수 있다.
  \item 정부종합청사는 행정재산 중 공용재산이다.
\end{choices}
\end{question}
\answer{1}
